\subsection{Five four-element corners with natural graph chains}

\begin{lemma}[Proved: the graph-chain version of the order-six obstruction]
\label{CP2:A-six-graph-chain-obstruction}
Let $M,N$ be finite monoids, with $|M|\le6$ and
$\widetilde H_*(BN;\mathbb Q)=0$, and let $S$ be a
nonempty left-$N$, right-$M$ biset.  Suppose the augmented
chain complex of $B(N,S)$ has a natural, right-$M$-equivariant
chain homotopy equivalence to the augmented cellular
complex of a finite graph $\Gamma$.  Require the $M$
action on that graph to send each parametrized edge at
unit speed to an edge, possibly with reversed orientation,
or to a vertex.  If the collage has bounded rational
homology, all its rational degree-two products vanish.
The comparison of chains need not come from a map of spaces.
\end{lemma}
\begin{proof}
Collapse the corner $BN$ in the actual two-block homotopy
pushout.  This gives a rational homology equivalence,
hence a rational cohomology-ring isomorphism, to
\[
 Y=\operatorname{hocofib}\bigl(B(N,S,M)\longrightarrow BM\bigr).
\]
Write $R=\mathbb QM$ and $L_*$ for a projective left
$R$ resolution of the trivial module.  The given natural
chain equivalence computes the augmented bridge projection
by the graph complex tensored with $L_*$.  In particular,
putting
\[
 C=[C_1(\Gamma;\mathbb Q)\longrightarrow
             \ker(C_0(\Gamma;\mathbb Q)\to\mathbb Q)],
\]
the reduced chains of $Y$ are quasi-isomorphic to
$(C\otimes_R^{\mathbf L}\mathbb Q)[1]$.  This follows
from the degreewise exact augmented sequence, since
the terms of $L_*$ are flat.

If $H^2(BM;\mathbb Q)=0$, every second class of $Y$
lifts from its actual suspension quotient, so its
products vanish.  Otherwise the small-minimal-ideal
argument in Theorem~\ref{CP2:A-baseline-six-graph} gives
a $2\times2$ rectangular minimal band $K$, with at most
two elements outside it.  Let $e=x_{00}$ and $e_1=x_{01}$.
If the two image graphs $\Gamma e,\Gamma e_1$ share a
vertex, that vertex is fixed by all of $M$, by the
literal band multiplication.  It supplies a right-$R$
chain section of the graph augmentation.  After tensoring
with $L_*$ and using the given augmented comparison,
the bridge projection is therefore injective on the
pullback of $H^2(BM;\mathbb Q)$.  The cofiber exact
sequence makes every second class of $Y$ restrict to
zero on $BM$.  Its lift from the actual suspension
quotient again has zero products.  Only this algebraic
section is needed; no graph-to-bridge map is asserted.

If the image graphs have disjoint vertex sets, they
also have disjoint edge sets.  The literal two-column
map $Ce\otimes eR\to C$ is injective in both degrees:
in degree zero both column copies separately have
augmentation zero and disjoint vertex supports; in
degree one their edge supports are disjoint.  Its
image is $C\mathbb QK$.  Thus
$D=C/(Ce\otimes eR)$ is a two-term complex over
$A=R/\mathbb QK$.  The subcomplex removed is right
$R$-projective term by term.  With $V$ the one-dimensional
row augmentation, the exact sequence
$0\to V\to Re\to\mathbb Q\to0$ gives
\[
 H_n(Y;\mathbb Q)
   =H_{n-2}(D\otimes_A^{\mathbf L}V),\qquad n\ge4.
\]
Here the change of rings follows from the left-projective
resolution $0\to\mathbb QK\to R\to A\to0$ and
$De=0$, exactly as in the cited graph theorem.
The algebra $A$ is semisimple or the dual-number algebra.
In the first case $V$ is projective and the degree-two
group on the right is zero.  In the second, $V=k$ and
$0\to k\to A\to k\to0$ gives
\[
 H_j(D\otimes_A^{\mathbf L}k)
   \cong H_{j-1}(D\otimes_A^{\mathbf L}k),\qquad j\ge3.
\]
Boundedness forces their degree-two group to be zero,
so $H_4(Y;\mathbb Q)=0$.  This proves the assertion
without using a multiplicative chain comparison.
\end{proof}

\begin{lemma}[Proved: explicit graph complexes for the five corners]
\label{CP2:A-four-corner-graph-chain}
Let $N$ be one of the following monoids:
the identity-adjoined right-zero band on two or three
elements, or a monoid on $\{1,e_0,e_1,a\}$ with
\[
 e_i e_j=e_j,\qquad ae_i=e_i,\qquad
 e_i a=e_{\tau(i)},
\]
where $(\tau,a^2)$ is one of
\[
 (\mathrm{id},1),\quad(\mathrm{id},a),\quad
 ((01),1),\quad(\mathrm{const}_0,a).
\]
Every finite left $N$-set has a natural augmented graph
chain model of the kind in the preceding lemma, for
every commuting right monoid action.
\end{lemma}
\begin{proof}
Put $T=\mathbb QN$.  The $e_i$ actors are retractions
$p_i$ with common image $U$; each fixes $U$ pointwise.
Thus $e_0\mathbb QS=\mathbb QU$.  The right-projective
module $e_0T$ has basis $e_i$, and the kernel of its
augmentation has basis $e_i-e_0$ for $i>0$.

For an identity-adjoined right-zero band of size $c$,
that kernel is the direct sum of $c-1$ copies of
$(1-e_0)T$.  Each copy maps $1-e_0$ to $e_i-e_0$.
Consequently this is an exact projective resolution of
the trivial right module of length one.  After tensoring
with $\mathbb QS$, it is the graph with vertex set $U$
and, for each $s\notin U$ and $i>0$, an edge from
$p_0s$ to $p_i s$.  A basis vector is $s-p_0s$ in
the $i$th copy; a commuting map sends it to the
corresponding basis vector or to zero, collapsing the
edge to its common endpoint in the latter case.

For the four remaining monoids the kernel is
$\mathbb Q(e_1-e_0)$.  The following table gives an
idempotent $q$ such that $qT$ is one-dimensional and
$q\mapsto e_1-e_0$ is a right-module isomorphism to
that kernel:
\[
\begin{array}{c|c|c}
\tau&a^2&q\\\hline
\mathrm{id}&1&(1+a)/2-e_0\\
\mathrm{id}&a&a-e_0\\
(01)&1&(1-a)/2\\
\mathrm{const}_0&a&1-a
\end{array}
\]
All claims in this table follow directly from the
displayed products: $q^2=q$, $qe_i=0$, and the right
action of $a$ on $q$ is respectively $+1,+1,-1,0$,
the same as on $e_1-e_0$.  We therefore have the
exact right-projective resolution
\[
 0\longrightarrow qT\longrightarrow e_0T
   \longrightarrow\mathbb Q\longrightarrow0.
\]
Write $A$ for the actor of $a$ on $S$.  In each case
the degree-zero term after tensoring is $\mathbb QU$.
The degree-one term $q\mathbb QS$ has the following
explicit graph bases.

For $(\mathrm{id},1)$, $A$ is an involution fixing
$U$ pointwise.  Use one edge for each $A$-orbit $O$
outside $U$, represented by
$|O|^{-1}\sum_{s\in O}s-p_0s$.  Both $p_i$ are
constant on $O$, so its boundary is $p_1s-p_0s$.
A commuting map sends an orbit average to the average
of its image orbit, and hence sends this basis vector
to another such vector or to zero.

For $(\mathrm{id},a)$, use the basis $t-p_0t$ for
$t\in AS\setminus U$, with boundary $p_1t-p_0t$.
Commuting maps preserve $AS$ and $U$, and again give
edge maps or collapses.

For $((01),1)$, use one oriented edge for each
two-element $A$-orbit, with basis vector $(s-As)/2$.
Its boundary is $p_1s-p_0s$, since $p_0A=p_1$ and
$p_1A=p_0$.  Choosing the other representative reverses
the orientation.  A commuting map sends this vector
to a basis vector, its negative, or zero.  In the
last case its image is $A$-fixed and its two endpoints
coincide, so the graph map collapses to that vertex.

For $(\mathrm{const}_0,a)$, use the basis $s-As$ for
$s\notin AS$.  Its boundary is $p_1s-p_0s$, because
$p_iA=p_0$.  A commuting map gives another such basis
vector or zero; in the latter case the common endpoint
again specifies the collapsed vertex.  These rules
define strict unit-speed graph actions in all cases.

Comparison of the displayed projective resolutions
with the ordinary right bar resolution gives an
augmented chain homotopy equivalence.  Its tensor with
$\mathbb QS$ is natural in the $N$-set $S$, so the
comparison is equivariant for every commuting monoid
action.  This establishes the precise chain-level
hypothesis, without asserting an unconstructed graph
equivalence of action spaces.
\end{proof}

\begin{theorem}[Proved: all but one four-element opposite corner]
\label{CP2:A-six-versus-four-semisimple-quotient}
Let $S$ be a nonempty left-$N$, right-$M$ biset with
$|M|\le6$ and $|N|\le4$.  Suppose $N$ is not isomorphic
to the monoid
\[
 N_{\mathrm{nil}}=\{1,e_0,e_1,a\},\qquad
 e_i e_j=e_j,\quad ae_i=e_i,\quad e_i a=e_0,
 \quad a^2=e_0.
\]
If the collage has bounded rational homology, all its
rational degree-two products vanish.  There is no
restriction on $|S|$.  The statement for the opposite
collage follows by taking opposites.
\end{theorem}
\begin{proof}
If the minimal ideal of $N$ has one column, averaging
its maximal group gives an idempotent $p\in\mathbb QN$
with $p\mathbb QN=\mathbb Qp$ and augmentation one.
Thus the trivial right module is projective.  The
actual orbit functor, replacing $N$ by the trivial
corner and $S$ by $N\backslash S$, is a rational
homology equivalence on the collage.  Its target is
an ordinary act cone over $M$, whose rational second
products vanish by Lemma~\ref{CP2:A-high-seven-act}.

Otherwise the minimal ideal has at least two columns.
The order-four budget forces a trivial maximal group,
one row, and two or three columns; the global identity
is outside it.  Three columns give the identity-adjoined
three-element right-zero band.  Two columns give the
three-element right-zero monoid, or one additional
element $a$ besides the global identity.  Necessarily
$ae_i=e_i$ and right multiplication by $a$ induces a
transformation $\tau$ of the two columns.  If $a^2=1$,
then $\tau$ is the identity or the transposition.  If
$a^2=a$, it is the identity or a constant.  If $a^2$
lies in the band, then $\tau$ is constant and its
value equals $a^2$, up to exchanging the columns.
This lists precisely the four cases of the preceding
lemma and the excluded nilpotent case.

Every remaining listed corner has contractible
classifying space, since any $e_i$ is a right zero
for the whole monoid.  Apply the preceding graph-chain
lemma and the order-six graph-chain obstruction.
\end{proof}
